% ============================================================
%  The GGM construction
% ============================================================

\section{Pseudorandom functions from PRGs: the GGM construction}

\Cref{thm:prf-cpa} is useful only if PRFs exist. Goldreich, Goldwasser and Micali (GGM) showed how to build a PRF from any PRG, so PRFs exist whenever one-way functions do.

\subsection{Input of one bit}

Start from a \textbf{length-doubling} PRG and split its output into two halves of $\lambda$ bits:
\[
  G : \{0,1\}^\lambda \to \{0,1\}^{2\lambda} \qquad\qquad G(s) = \big(G_0(s),\ G_1(s)\big)
\]

\begin{framedprop}[label={prop:ggm-one-bit}]{A PRF with domain $\{0,1\}$}
  The family $F_k(x) = G_x(k)$, for $k \in \{0,1\}^\lambda$ and $x \in \{0,1\}$, is a PRF.
\end{framedprop}

The key is the seed of the PRG, and the input bit says \textbf{which half} to output: $F_k(0) = G_0(k)$ (left half) and $F_k(1) = G_1(k)$ (right half).

\textbf{Why it works.} A function with only two inputs is a table with \textbf{two rows}. A truly random such function is just two independent random values, that is, $2\lambda$ random bits. Here the whole table is $\big(F_k(0), F_k(1)\big) = G(k)$, which by the security of the PRG looks like $2\lambda$ random bits. So whatever the adversary asks, its view in $\mathrm{REAL}$ is indistinguishable from its view in $\mathrm{RAND}$.

\subsection{Input of $n$ bits: the tree}

One input bit is not enough. The idea of GGM is to repeat the choice ``left or right'' \textbf{once for every bit of the input}, building a binary tree:
\begin{itemize}
  \item the \textbf{root} is the key $k$;
  \item every node $s$ has two children: $G_0(s)$ on the left and $G_1(s)$ on the right;
  \item to compute $F_k(x_1 x_2 \cdots x_n)$, start from the root and go down: at level $i$ go \textbf{left if $x_i = 0$} and \textbf{right if $x_i = 1$};
  \item the output is the \textbf{leaf} reached.
\end{itemize}

\begin{frameddefn}[label={def:ggm}]{GGM pseudorandom function}
  Let $G(s) = (G_0(s), G_1(s))$ be a length-doubling PRG. For $k \in \{0,1\}^\lambda$ and $x = x_1 \cdots x_n \in \{0,1\}^n$:
  \[
    F_k(x_1 x_2 \cdots x_n) = G_{x_n}\Big(\cdots G_{x_2}\big(G_{x_1}(k)\big)\cdots\Big)
  \]
\end{frameddefn}

\begin{figure}[H]
  \centering
  \begin{tikzpicture}[
      level 1/.style={sibling distance=62mm, level distance=14mm},
      level 2/.style={sibling distance=31mm, level distance=14mm},
      every node/.style={font=\small},
      path/.style={very thick, BrickRed},
      edge from parent/.style={draw, thin}]
    \node {$k$}
      child {node {$G_0(k)$}
        child {node {$G_0(G_0(k))$}}
        child {node {$G_1(G_0(k))$}}
        edge from parent node[left, lbl] {$0$}}
      child {node[text=BrickRed] {$G_1(k)$}
        child {node[text=BrickRed] {$\mathbf{G_0(G_1(k))}$} edge from parent[path] node[left, lbl] {$0$}}
        child {node {$G_1(G_1(k))$}}
        edge from parent[path] node[right, lbl] {$1$}};
  \end{tikzpicture}
  \caption{The GGM tree for $n = 2$. The four leaves are the four values of $F_k$; the highlighted path computes $F_k(10) = G_0(G_1(k))$.}
\end{figure}

\begin{itemize}
  \item \textbf{The tree is the giant table} $\to$ it has $2^n$ leaves, one for every input.
  \item \textbf{It is efficient} $\to$ to compute $F_k(x)$ we only walk \textbf{one path}, so $n$ evaluations of $G$ are enough; the rest of the tree is never computed.
\end{itemize}

\begin{framedthm}[label={thm:ggm}]{Goldreich, Goldwasser, Micali}
  If $G$ is a length-doubling PRG, then the family of \Cref{def:ggm} is a PRF.
\end{framedthm}

\textbf{Proof idea.} A hybrid argument over the levels of the tree. In the hybrid $H_i$ the nodes of level $i$ are \textbf{truly random} (sampled lazily, only where the adversary's queries pass) and the levels below are computed with $G$. $H_0$ is the real PRF (only the root is random), and $H_n$ is a truly random function (all leaves are random). Two adjacent hybrids differ in replacing $G(s)$ by random strings on the nodes of one level that the adversary touches: these are at most polynomially many independent evaluations of $G$ on random seeds, which are indistinguishable from random by the security of the PRG (with a further hybrid over the evaluations). With $n$ levels and polynomially many queries, the total loss is negligible.

\begin{framedcor}{From one-way functions to CPA security}
  \[
    \text{OWF} \;\xrightarrow{\text{HILL}}\; \text{PRG} \;\xrightarrow{\text{GGM}}\; \text{PRF} \;\xrightarrow{\text{\Cref{thm:prf-cpa}}}\; \text{CPA-secure encryption}
  \]
  If one-way functions exist, then CPA-secure symmetric encryption exists.
\end{framedcor}
